% GENERATED FILE. Edit canonical lessons, then run npm run generate:books.
% canonical-source-hash: fnv1a64:4df622d099168245
% canonical-lessons: MR-C284-rag, MR-C284-bhiti, MR-C284-rag-ala, MR-C284-dhyey, MR-C284-vhaycha-ahe, MR-R284-feelings-and-hopes

\chapter{Feelings and Hopes}
\label{ch:mr-feelings-and-hopes}
\hlchaptermodality{284}

\begin{chapteropening}
\textbf{By the end of this chapter:} \emph{I can say what I felt (malā rāg ālā, bhītī vāṭlī), name a goal (dhyey), and say what I want to do and become (vhāycaṁ āhe).}

Complete MR-R284-feelings-and-hopes: say what you felt and what you want to become.
\end{chapteropening}

\section[\texorpdfstring{\mr{राग} (rāg)}{राग (rāg)}]{\mr{राग} (rāg) — anger}

\label{lesson:MR-C284-rag}



\begin{quote}
Before the new one: say \textquotedblleft{}my bag is lost\textquotedblright{} and \textquotedblleft{}I missed my train\textquotedblright{}. (\emph{mājhī piśvī haravlī}, \emph{mājhī gāḍī cuklī}.) Then \textquotedblleft{}to get angry\textquotedblright{}. (\emph{rāgāvṇe}.)
\end{quote}

\subsection*{You'll want to know: \mr{राग}}
\textbf{\mr{राग}} — \emph{rāg} — \textquotedblleft{}anger\textquotedblright{}.

\begin{quote}
\textbf{\mr{मला} \mr{खूप} \mr{राग} \mr{आला}.}
\end{quote}

\begin{quote}
\emph{malā khūp rāg ālā.} — \textquotedblleft{}I got very angry.\textquotedblright{}
\end{quote}

\textbf{\mr{रागावणे}}, \textquotedblleft{}to get angry\textquotedblright{}, is built on this noun. Anger is something that comes to you: \textquotedblleft{}to me anger came\textquotedblright{}.

\begin{cousinweb}
\textbf{\mr{राग}} (\emph{rāg}) is a tatsama: Sanskrit \emph{rāga}, \textquotedblleft{}passion, love\textquotedblright{} (CDIAL 10674); CDIAL lists no Marathi form. Marathi, like Bengali, uses it for anger; how the sense moved there is not settled. Sanskrit \emph{rāga} also names a melody, the \emph{rāga} of Indian music (Monier-Williams).
\end{cousinweb}

\subsection*{Your turn}
Say these aloud:

\begin{itemize}
\raggedright
  \item \emph{rāg}
  \item \textquotedblleft{}I got very angry\textquotedblright{} — \emph{malā khūp rāg ālā}
  \item \emph{rāg}, then \emph{rāgāvṇe}, the verb built on it
\end{itemize}

\begin{tcolorbox}[breakable,colback=teal!4,colframe=teal!35!black,title={Before you move on}]
What does \textbf{\mr{राग}} mean? (\textquotedblleft{}Anger\textquotedblright{}.) Which verb is built on it? (\textbf{\mr{रागावणे}}.)
\end{tcolorbox}
\section[\texorpdfstring{\mr{भीती} (bhītī)}{भीती (bhītī)}]{\mr{भीती} (bhītī) — fear}

\label{lesson:MR-C284-bhiti}



\begin{quote}
Before the new one: say \textquotedblleft{}anger\textquotedblright{} and \textquotedblleft{}to be afraid\textquotedblright{}. (\emph{rāg}, \emph{ghābarṇe}.)
\end{quote}

\subsection*{You'll want to know: \mr{भीती}}
\textbf{\mr{भीती}} — \emph{bhītī} — \textquotedblleft{}fear\textquotedblright{}.

\begin{quote}
\textbf{\mr{मला} \mr{सापाची} \mr{भीती} \mr{वाटते}.}
\end{quote}

\begin{quote}
\emph{malā sāpācī bhītī vāṭte.} — \textquotedblleft{}I am afraid of snakes.\textquotedblright{}
\end{quote}

What you fear takes \textbf{-\mr{ची}}: \textbf{\mr{सापाची} \mr{भीती}}, \textquotedblleft{}fear of snakes\textquotedblright{}. The feeling \textbf{\mr{वाटते}}, \textquotedblleft{}is felt\textquotedblright{}, and it is feminine, so the verb ends in \textbf{-\mr{ते}}.

\begin{cousinweb}
\textbf{\mr{भीती}} (\emph{bhītī}) is a tatsama: Sanskrit \emph{bhīti}, \textquotedblleft{}fear\textquotedblright{} (CDIAL 9512), from the root \emph{bhī}, \textquotedblleft{}to be afraid\textquotedblright{}. Marathi writes it with a long final \emph{ī}.
\end{cousinweb}

\subsection*{Your turn}
Say these aloud:

\begin{itemize}
\raggedright
  \item \emph{bhītī}
  \item \textquotedblleft{}I am afraid of snakes\textquotedblright{} — \emph{malā sāpācī bhītī vāṭte}
  \item \emph{bhītī}, once more, and \emph{ghābarṇe}, \textquotedblleft{}to be afraid\textquotedblright{}
\end{itemize}

\begin{tcolorbox}[breakable,colback=teal!4,colframe=teal!35!black,title={Before you move on}]
What does \textbf{\mr{भीती}} mean? (\textquotedblleft{}Fear\textquotedblright{}.) \textquotedblleft{}I am afraid of snakes\textquotedblright{}? (\textbf{\mr{मला} \mr{सापाची} \mr{भीती} \mr{वाटते}}.)
\end{tcolorbox}
\section[\texorpdfstring{\mr{मला} \mr{राग} \mr{आला} (malā rāg ālā)}{मला राग आला (malā rāg ālā)}]{\mr{मला} \mr{राग} \mr{आला} (malā rāg ālā) — feelings that come to you}

\label{lesson:MR-C284-rag-ala}



\begin{quote}
Say \textquotedblleft{}I got very angry\textquotedblright{} and \textquotedblleft{}I am afraid of snakes\textquotedblright{}. (\emph{malā khūp rāg ālā}, \emph{malā sāpācī bhītī vāṭte}.) Then the first line you learned: \emph{bheṭūn ānand jhālā}.
\end{quote}

\subsection*{You'll want to know: the feeling is the subject}
\begin{quote}
\textbf{\mr{मला} \mr{राग} \mr{आला}.}
\end{quote}

\begin{quote}
\emph{malā rāg ālā.} — \textquotedblleft{}I got angry.\textquotedblright{}
\end{quote}

The one who feels stands in the dative, as with \textbf{\mr{आवडणे}} and \textbf{\mr{वाटणे}}. The feeling is the subject, and the past agrees with it: \textbf{\mr{राग} \mr{आला}} (masculine), \textbf{\mr{भीती} \mr{वाटली}} (feminine), \textbf{\mr{कंटाळा} \mr{आला}} (\emph{kaṇṭāḷā ālā}), \textquotedblleft{}I got bored\textquotedblright{}. Ask it the same way: \textbf{\mr{तुम्हाला} \mr{राग} \mr{आला} \mr{का}?} (\emph{tumhālā rāg ālā kā?})

Two feeling words you know are adjectives instead, and they agree with the person: \textbf{\mr{ती} \mr{आनंदी} \mr{आहे}} (\emph{tī ānandī āhe}), \textquotedblleft{}she is happy\textquotedblright{}; \textbf{\mr{तो} \mr{दुःखी} \mr{होता}} (\emph{to duḥkhī hotā}), \textquotedblleft{}he was sad\textquotedblright{}.

\subsection*{You'll want to know: which verb each feeling takes}
Three verbs carry feelings, and each feeling has its own:

\begin{itemize}
\raggedright
  \item \textbf{\mr{आला}}, \textquotedblleft{}came\textquotedblright{}: \textbf{\mr{राग} \mr{आला}}, \textbf{\mr{कंटाळा} \mr{आला}}
  \item \textbf{\mr{वाटलं}}, \textquotedblleft{}was felt\textquotedblright{}: \textbf{\mr{भीती} \mr{वाटली}}, \textbf{\mr{लाज} \mr{वाटली}} (\emph{lāj vāṭlī}), \textquotedblleft{}I felt ashamed\textquotedblright{}
  \item \textbf{\mr{झालं}}, \textquotedblleft{}came about\textquotedblright{}: \textbf{\mr{आनंद} \mr{झाला}}, \textbf{\mr{दुःख} \mr{झालं}} (\emph{duḥkh jhālaṁ}), \textquotedblleft{}I was sad\textquotedblright{}
\end{itemize}

Learn the feeling with its verb, as you learned \textbf{\mr{भेटून} \mr{आनंद} \mr{झाला}} whole.

\subsection*{Your turn}
Say these aloud:

\begin{itemize}
\raggedright
  \item \textquotedblleft{}I got angry\textquotedblright{} — \emph{malā rāg ālā}
  \item \textquotedblleft{}She felt afraid\textquotedblright{} — \emph{tilā bhītī vāṭlī}
  \item \textquotedblleft{}We were glad\textquotedblright{} — \emph{āmhālā ānand jhālā}
  \item \textquotedblleft{}He got bored\textquotedblright{} — \emph{tyālā kaṇṭāḷā ālā}
\end{itemize}

\textbf{\mr{आम्हाला}} (\emph{āmhālā}, \textquotedblleft{}to us\textquotedblright{}).

\begin{tcolorbox}[breakable,colback=teal!4,colframe=teal!35!black,title={Before you move on}]
\textquotedblleft{}I got angry\textquotedblright{}? (\textbf{\mr{मला} \mr{राग} \mr{आला}}.) \textquotedblleft{}I felt ashamed\textquotedblright{}? (\textbf{\mr{मला} \mr{लाज} \mr{वाटली}}.) \textquotedblleft{}We were glad\textquotedblright{}? (\textbf{\mr{आम्हाला} \mr{आनंद} \mr{झाला}}.) What does the verb agree with? (The feeling.)
\end{tcolorbox}
\section[\texorpdfstring{\mr{ध्येय} (dhyey)}{ध्येय (dhyey)}]{\mr{ध्येय} (dhyey) — a goal}

\label{lesson:MR-C284-dhyey}



\begin{quote}
Before the new one: say \textquotedblleft{}a wish\textquotedblright{} and \textquotedblleft{}a dream\textquotedblright{}. (\emph{icchā}, \emph{svapna}.) Then \textquotedblleft{}we were glad\textquotedblright{}. (\emph{āmhālā ānand jhālā}.)
\end{quote}

\subsection*{You'll want to know: \mr{ध्येय}}
\textbf{\mr{ध्येय}} — \emph{dhyey} — \textquotedblleft{}a goal, an aim\textquotedblright{}.

\begin{quote}
\textbf{\mr{माझं} \mr{ध्येय} \mr{मोठं} \mr{आहे}.}
\end{quote}

\begin{quote}
\emph{mājhaṁ dhyey moṭhaṁ āhe.} — \textquotedblleft{}My goal is a big one.\textquotedblright{}
\end{quote}

\textbf{\mr{ध्येय}} is neuter: \textbf{\mr{माझं} \mr{ध्येय}}. A wish can stay a wish; a goal is what you work toward.

\begin{cousinweb}
\textbf{\mr{ध्येय}} (\emph{dhyey}) is a tatsama: Sanskrit \emph{dhyeya}, \textquotedblleft{}to be meditated on, to be pondered\textquotedblright{} (Monier-Williams), from the root of \emph{dhyāna}, \textquotedblleft{}meditation\textquotedblright{}. Marathi uses it for the aim one keeps the mind fixed on.
\end{cousinweb}

\subsection*{Your turn}
Say these aloud:

\begin{itemize}
\raggedright
  \item \emph{dhyey}
  \item \textquotedblleft{}My goal is a big one\textquotedblright{} — \emph{mājhaṁ dhyey moṭhaṁ āhe}
  \item a wish, a dream and a goal: \emph{icchā}, \emph{svapna}, \emph{dhyey}
\end{itemize}

\begin{tcolorbox}[breakable,colback=teal!4,colframe=teal!35!black,title={Before you move on}]
What does \textbf{\mr{ध्येय}} mean? (\textquotedblleft{}A goal, an aim\textquotedblright{}.) Is it masculine, feminine or neuter? (Neuter: \textbf{\mr{माझं} \mr{ध्येय}}.)
\end{tcolorbox}
\section[\texorpdfstring{\mr{मला} \mr{शिक्षक} \mr{व्हायचं} \mr{आहे} (malā śikṣak …)}{मला शिक्षक व्हायचं आहे (malā śikṣak …)}]{\mr{मला} \mr{शिक्षक} \mr{व्हायचं} \mr{आहे} (malā śikṣak vhāycaṁ āhe) — wanting to do, wanting to be}

\label{lesson:MR-C284-vhaycha-ahe}



\begin{quote}
Say \textquotedblleft{}my goal is a big one\textquotedblright{}. (\emph{mājhaṁ dhyey moṭhaṁ āhe}.) Then \textquotedblleft{}how do I get to the station?\textquotedblright{}, with its \textbf{\mr{जायचं}}. (\emph{sṭeśanlā kasaṁ jāycaṁ?})
\end{quote}

\subsection*{You'll want to know: \mr{मला} … -\mr{ायचं} \mr{आहे}, 'I want to'}
\begin{quote}
\textbf{\mr{मला} \mr{मराठीत} \mr{गोष्ट} \mr{लिहायची} \mr{आहे}.}
\end{quote}

\begin{quote}
\emph{malā marāṭhīt goṣṭa lihāycī āhe.} — \textquotedblleft{}I want to write a story in Marathi.\textquotedblright{}
\end{quote}

The \textbf{-\mr{ायच}-} form, with \textbf{\mr{आहे}} and the dative: \textbf{\mr{मला} \mr{जायचं} \mr{आहे}} (\emph{malā jāycaṁ āhe}), \textquotedblleft{}I want to go\textquotedblright{}. With an object it agrees with the thing, as the past does: \textbf{\mr{गोष्ट} \mr{लिहायची} \mr{आहे}}, \textbf{\mr{पुस्तक} \mr{वाचायचं} \mr{आहे}} (\emph{pustak vācāycaṁ āhe}). And the \textquotedblleft{}not\textquotedblright{}: \textbf{\mr{मला} \mr{जायचं} \mr{नाही}} (\emph{malā jāycaṁ nāhī}), \textquotedblleft{}I don't want to go\textquotedblright{}.

\subsection*{You'll want to know: \mr{व्हायचं}, 'to become'}
\begin{quote}
\textbf{\mr{मला} \mr{शिक्षक} \mr{व्हायचं} \mr{आहे}.}
\end{quote}

\begin{quote}
\emph{malā śikṣak vhāycaṁ āhe.} — \textquotedblleft{}I want to become a teacher.\textquotedblright{}
\end{quote}

\textbf{\mr{व्हायचं}} is \textbf{\mr{होणे}}, \textquotedblleft{}to become\textquotedblright{}, in the same form. The wish words you know can follow it: \textbf{\mr{ते} \mr{माझं} \mr{स्वप्न} \mr{आहे}} (\emph{te mājhaṁ svapna āhe}), \textquotedblleft{}that is my dream\textquotedblright{}; \textbf{\mr{ते} \mr{माझं} \mr{ध्येय} \mr{आहे}} (\emph{te mājhaṁ dhyey āhe}), \textquotedblleft{}that is my goal\textquotedblright{}.

\subsection*{Your turn}
Say these aloud:

\begin{itemize}
\raggedright
  \item \textquotedblleft{}I want to become a teacher\textquotedblright{} — \emph{malā śikṣak vhāycaṁ āhe}
  \item \textquotedblleft{}I want to write a story in Marathi\textquotedblright{} — \emph{malā marāṭhīt goṣṭa lihāycī āhe}
  \item \textquotedblleft{}I don't want to go\textquotedblright{} — \emph{malā jāycaṁ nāhī}
\end{itemize}

\begin{tcolorbox}[breakable,colback=teal!4,colframe=teal!35!black,title={Before you move on}]
\textquotedblleft{}I want to go\textquotedblright{}? (\textbf{\mr{मला} \mr{जायचं} \mr{आहे}}.) \textquotedblleft{}I want to become a teacher\textquotedblright{}? (\textbf{\mr{मला} \mr{शिक्षक} \mr{व्हायचं} \mr{आहे}}.) \textquotedblleft{}I don't want to go\textquotedblright{}? (\textbf{\mr{मला} \mr{जायचं} \mr{नाही}}.)
\end{tcolorbox}
\section[(review)]{Review — feelings and hopes}

\label{lesson:MR-R284-feelings-and-hopes}



\begin{quote}
Ask for one ticket to Pune and when the train leaves. (\emph{puṇyācaṁ ek tikīṭ dyā. gāḍī kitī vājtā nighte?}) Say \textquotedblleft{}I bought a ticket, but I missed my train\textquotedblright{}. (\emph{mī tikīṭ kāḍhlaṁ, paṇ mājhī gāḍī cuklī}.) Then \textbf{\mr{ऋतू}} and \textbf{\mr{ज्ञान}}: two words, two letters, \textbf{\mr{ऋ}} and \textbf{\mr{ञ}}.
\end{quote}

\subsection*{Your turn}
Say these aloud:

\begin{itemize}
\raggedright
  \item I missed my train, so I got angry.
  \item She is afraid of snakes.
  \item We were glad. He got bored.
  \item I want to become a teacher. That is my goal.
  \item one thing you want to do this year, with \textbf{-\mr{ायचं} \mr{आहे}}
\end{itemize}

\begin{tcolorbox}[breakable,colback=teal!4,colframe=teal!35!black,title={Before you move on}]
Anger? (\textbf{\mr{राग}}.) Fear? (\textbf{\mr{भीती}}.) \textquotedblleft{}I got angry\textquotedblright{}? (\textbf{\mr{मला} \mr{राग} \mr{आला}}.) A goal? (\textbf{\mr{ध्येय}}.) \textquotedblleft{}I want to become\textquotedblright{}? (\textbf{\mr{मला} … \mr{व्हायचं} \mr{आहे}}.) Which letter starts \textbf{\mr{ऋतू}}? (\textbf{\mr{ऋ}}.)
\end{tcolorbox}
